\documentclass[11pt]{article}
\usepackage{amsmath,amssymb,amsthm,hyperref}
\newtheorem{theorem}{Theorem}
\newtheorem{lemma}[theorem]{Lemma}
\newcommand{\A}{\mathcal A}
\title{Subset divisibility completely characterizes signed fair signatures}
\author{Research proof; two independent mathematical reviews completed}
\date{30 September 2026}
\begin{document}
\maketitle

Let $\A_{\mathbb Q}$ be the associative algebra with basis the empty word
and all words of distinct letters from $X_1,\ldots,X_n$. Multiplication
is concatenation when the two letter sets are disjoint, and zero
otherwise. Equivalently, it is the free associative algebra modulo the
ideal generated by all words containing a repeated letter.
Let $\A_{\mathbb Z}$ denote its integer-coefficient subring and let
$\mathfrak m$ denote the positive-degree ideal. This ideal is nilpotent,
so exponential and logarithm are finite polynomials over $\mathbb Q$.

A signed word is a product of the elements
$1+X_i=\exp(X_i)$ and their inverses $1-X_i=\exp(-X_i)$.
Negative letters are purely algebraic: they are not legal die faces.

\paragraph{Classical algebraic setting.}
This is the reduced free associative algebra used in Milnor's Magnus
expansion of the reduced free group. The group generated by $1+X_i$
identifies with that reduced free group by
\cite[Theorem~1.12 and Corollary~1.13]{darne}.
The integral fixed-last-letter Lie basis used below is also classical;
see \cite[Lemma~6.15]{darne}. The elementary proofs are retained to
make the application to the fair target self-contained.
For positive words, equality of these reduced signatures is exactly
strong $M$-equivalence in \cite[Theorem~3.3]{teh}.
No novelty is claimed for these algebraic tools.

\begin{theorem}\label{main}
For integers $c_1,\ldots,c_n$, the fair signature
\[
 T_c=\exp\left(\sum_i c_iX_i\right)
\]
is represented by a signed word if and only if
\[
 |S|!\ \bigm|\ \prod_{i\in S}c_i
 \quad\text{for every nonempty }S\subseteq[n].
\]
Thus integral membership in the reduced nilpotent group supplies no
arithmetic obstruction beyond the subset-factorial divisibility already
known for ordinary fair words. Any stronger obstruction for positive words
must use positivity or some other semigroup constraint.
\end{theorem}

\begin{lemma}[An integral multilinear Lie basis]\label{basis}
Fix a nonempty set of letters $S$ and a distinguished $a\in S$.
For each ordering $\pi$ of $S\setminus\{a\}$, put
\[
 B_\pi=[X_{\pi_1},[X_{\pi_2},\ldots,
                     [X_{\pi_{|S|-1}},X_a]\ldots]].
\]
For $S=\{a\}$ use $B_\varnothing=X_a$.
These form a basis of the multilinear Lie component on $S$ over
$\mathbb Q$. In this basis, the coefficient of $B_\pi$ in a Lie element
$L$ is exactly the associative-word coefficient of
$X_{\pi_1}\cdots X_{\pi_{|S|-1}}X_a$ in $L$.
In particular, an integral-coefficient multilinear Lie element has
integer coordinates in this basis.
\end{lemma}
\begin{proof}
Antisymmetry and the Jacobi identity rewrite every multilinear Lie
bracket as a linear combination of brackets right-nested toward the
fixed letter $a$. One can see this inductively from
$[[U,V],W]=[U,[V,W]]-[V,[U,W]]$, always moving the bracket containing
$a$ to the right. Thus the displayed brackets span.
When $B_\pi$ is expanded associatively, its only word ending in $X_a$
is $X_{\pi_1}\cdots X_{\pi_{|S|-1}}X_a$, with coefficient one.
The ending-$a$ coefficients therefore prove linear independence and
the asserted coordinate formula.
\end{proof}

\begin{proof}[Proof of Theorem~\ref{main}]
The coefficient of every word whose letter set is $S$ in $T_c$ is
$\prod_{i\in S}c_i/|S|!$. Every signed word has integer coefficients,
so the divisibility conditions are necessary.

For sufficiency, the conditions say $T_c\in\A_{\mathbb Z}$.
We construct a signed word $g$ equal to $T_c$ by correcting degrees
successively. Suppose $g^{-1}T_c$ agrees with $1$ in degrees below $k$.
Because $g$ and $g^{-1}$ have integer coefficients, the homogeneous
part $L_k$ of degree $k$ of $g^{-1}T_c-1$ has integer coefficients.
It is also a Lie element: $\log g$ is a Lie polynomial by the finite
Baker--Campbell--Hausdorff formula, $\log T_c=\sum c_iX_i$ is a Lie
polynomial, and hence $\log(g^{-1}T_c)$ is one too. Its lowest
nonzero homogeneous part is $L_k$.

Decompose $L_k$ by its letter sets. Lemma~\ref{basis} expresses each
component as an integer combination of the $B_\pi$.
Each factor $1+B_\pi$ is itself represented by a signed word, obtained
by the identically nested group commutator. In fact if $U,V$ have
disjoint fixed supports and positive degrees, then
\[
 (1+U)(1+V)(1-U)(1-V)=1+[U,V]
\]
whenever $U^2=V^2=UVU=VUV=0$, as holds for such fixed-support elements
in this algebra. Induction gives the exact assertion for $B_\pi$.
Also $(1+B_\pi)^b=1+bB_\pi$ for every integer $b$, because
$B_\pi^2=0$.

Multiply $g$ on the right by the product of these factors with their
integer coordinates. This product is $1+L_k$ modulo degrees greater
than $k$, so the new $g^{-1}T_c$ agrees with $1$ through degree $k$.
Starting with $g=1$ and continuing through degree $n$ constructs the
required signed word.
\end{proof}

\paragraph{Equal counts.}
For $c_i=M>0$, the divisibility conditions are equivalent to every prime
$p\le n$ dividing $M$. Thus the primorial $M=\prod_{p\le n}p$ always
has an integral signed fair signature. In particular $n=4,M=6$ and
$n=5,M=30$ pass all reduced-group integrality tests. This does not assert
that ordinary fair dice with these parameters exist.

\paragraph{Exact finite checks.}
The companion script \texttt{check\_signed\_signatures.py} performs the
entire degree-correction algorithm in the integer algebra, verifies the
homogeneous Lie-basis identities at every step, and compares all final
coefficients to $T_c$. It records compact commutator coordinates rather
than expanding the possibly long signed word.
\begin{thebibliography}{9}
\bibitem{darne} Jacques Darn\'e,
\emph{Milnor invariants of braids and welded braids up to homotopy},
Algebraic \& Geometric Topology 24 (2024), 1277--1320,
\href{https://doi.org/10.2140/agt.2024.24.1277}{doi:10.2140/agt.2024.24.1277}.
The referenced theorem numbering was checked in the author's
\href{https://arxiv.org/abs/1904.10677}{arXiv:1904.10677v2}.
\bibitem{teh} Wen Chean Teh,
\emph{Parikh matrices and strong $M$-equivalence},
International Journal of Foundations of Computer Science 27 (2016),
545--556,
\href{https://doi.org/10.1142/S0129054116500155}{doi:10.1142/S0129054116500155};
\href{https://arxiv.org/abs/1501.07354}{arXiv:1501.07354}.
\end{thebibliography}
\end{document}
